\pageTitre{AVANT-PROPOS}

La formation professionnelle, telle qu'elle est organisée au Cameroun par le Ministère de
l'Emploi et de la Formation professionnelle (MINEFOP), repose sur l'approche par compétences.
Elle associe l'acquisition de connaissances théoriques, l'apprentissage des gestes
professionnels en atelier et une immersion en entreprise. Le stage de fin de formation en est
l'aboutissement : il confronte l'apprenant aux réalités d'un chantier, à ses contraintes de
délais, de qualité et de sécurité, et à la vie d'une équipe de travail.

C'est dans ce cadre que nous avons effectué un stage de \InfoDuree{} au sein de
\InfoEntreprise, du \InfoDateDebut{} au \InfoDateFin. Affecté à la division plomberie sur le
chantier \InfoChantier, nous avons participé à la pose des réseaux de distribution d'eau froide
et d'eau chaude sanitaire, à l'installation des organes de régulation de pression et à
l'isolation des canalisations.

Bien que notre spécialité soit le froid et la climatisation, ce stage nous a montré à quel
point ces métiers sont liés aux fluides du bâtiment : mêmes matériaux, mêmes techniques
d'assemblage, mêmes exigences d'étanchéité et d'isolation thermique. Le présent rapport rend
compte de ce que nous avons observé, appris et réalisé. Il est également une étape de
réflexion sur notre projet professionnel.
